% Standalone analytic fragment. No numerical coverage is asserted here.
\begin{lemma}[Exact pricing on physical source polygons]
\label{lem:physical-source-polygon-price}
Fix an amplitude interval $[u_-,u_+]$, a source interval $[r_-,r_+]$,
and a cumulant interval. Suppose that, throughout the cumulant interval
and for every allowed shape, proved scalar inequalities imply
\[
 a_j u+r\le C_j,\qquad a_j>0\quad(1\le j\le m).
\]
Let $K$ be the intersection of these half-planes with the amplitude--source
rectangle, after any separately proved coordinate restrictions.
For rational inputs its vertices and the maximum on $K$ of
\[
 F(u,r)=u^2+L_u u+g r^2+L_r r+C
\]
are computable using rational arithmetic. For a nonempty $K$, it suffices
to check its vertices and, on each edge, the stationary point in the edge
interior when the restricted quadratic is concave. This includes segments
and singletons. An empty intersection contains no physical source.

For a native certificate interpolated in $r$, first split $K$ at every
source knot. On each resulting piece apply this rule to each of the four
cumulant--shape corner bounds, and take the maximum of the resulting
bounds. A fixed nonnegative contact multiplier may be included by changing
only $L_u,L_r,C$.
\end{lemma}
\begin{proof}
Every defining half-plane contains every physical source in the nominal
rectangle, so intersection and coordinate clipping preserve all physical
points. Rational line intersections give all vertices, retaining points
on the boundaries. If a full-dimensional polygon had a maximizing point
in its interior, sufficiently small horizontal displacements would stay
in the polygon. Their two objective values have average
$F(u,r)+\delta^2$, contradicting maximality. Thus its maximum lies on its
boundary. On the edge $(u,r)=(u_0,r_0)+t(\Delta u,\Delta r)$,
$0\le t\le1$, the quadratic coefficient is
$(\Delta u)^2+g(\Delta r)^2$. Endpoints suffice when this coefficient is
nonnegative; otherwise also check the unique stationary point if it lies
in $(0,1)$. All these operations are rational. The same calculation covers
the lower-dimensional cases.

A native source chord is valid only between successive saved knots.
Splitting at every knot therefore retains its validity on each closed
piece; shared boundaries have both valid endpoint bounds. For each fixed
cumulant--shape corner the chord changes only the linear and constant
coefficients in $r$. The already proved corner inequalities and positive
interpolation uniformly bound the whole cumulant--shape cell. Taking the
maximum over the four corner prices and all source pieces gives a uniform
bound. The scalar row chosen at a source endpoint is fixed throughout its
cumulant--shape cell; different endpoints may use different globally
valid scalar rows.

For a contact with parameters $a,B$, put $d^2=6B^2$ and let $A$ be its
uniform scalar price. The nonnegative multiplier $\lambda$ adds
\[
 -\lambda a u-\lambda r+
 \lambda\left(A+\frac{\beta d}{b}\right)
 +\frac{\lambda^2d^2}{4b}.
\]
Here $\beta$ is the normalized native Hermite coefficient; for a raw
Hermite R3 profile the product $\beta d$ is $6\beta_{\rm raw}B$.
The coefficient of $u^2$ remains one, so the same boundary computation
applies. The multiplier and scalar contact are fixed on the whole
cell. Selecting any one of several proved upper bounds on each cell
preserves validity.
\end{proof}

\begin{remark}[Uniform source planes]
For the accepted even sextic scalar rows, the moment price is affine in
shape and convex in cumulant: the sextic coefficient is nonnegative and
the coefficient multiplying $\kappa^{3/2}$ is positive throughout the
shape interval. Hence the largest of the four endpoint prices is a
uniform $C_j$. A single fixed scalar row is used for each plane over the
whole cumulant interval. Taking the smallest uniform price among several
such rows is valid; taking separate minima at the four corners would
require a separate justification.
\end{remark}
